% AUTO-GENERATED by research/publication.py — sources: results/cyber_range/experiments/latest.json:suites.defense
\begin{table}[t]
\centering
\small
\caption{Cyber Range: Defense Comparison}
\label{tab:range_def}
\begin{tabular}{lcccc}
\toprule
\textbf{Defense} & \textbf{Attack success after defense} & \textbf{Availability during attack (mean / min)} & \textbf{Legitimate traffic blocked} & \textbf{Recovery after defense (s)} \\
\midrule
Block source & 0.0\% & 99.8\% / 96.4\% & 0.0\% & 2.6 \\
Rate limit & 19.7\% & 99.6\% / 95.9\% & 0.0\% & 3.7 \\
Isolate target & 0.0\% & 10.9\% / 0.0\% & 27.4\% & 52.4 \\
Monitor only & --- & 89.7\% / 84.7\% & 0.0\% & --- \\
Increase inspection & 100.0\% & 89.7\% / 84.7\% & 0.0\% & 47.4 \\
\bottomrule
\end{tabular}
\\[2pt]\parbox{0.95\linewidth}{\scriptsize Mean over 9 scenarios, 10 seeds each, automatic response. Means hide that most scenarios are cheap for the server; with monitoring only, availability during the attack averaged 66\% (DOS; worst second 37\%) and 41\% (DDOS; worst second 25\%). Availability is modelled (capacity / load), not measured on a real service; isolation is lifted 5 s after the attack stops, so its recovery time reflects scenario length. '---' = not applicable.}
\end{table}
